\documentclass[12pt]{article}

% Packages
\usepackage[a4paper, margin=1in]{geometry} % Adjust margins
\usepackage{setspace} % For line spacing
\usepackage{tocloft} % Optional: better spacing for TOC
\usepackage[backend=biber,style=apa]{biblatex} % Load biblatex with APA style
\addbibresource{references.bib} % Link to your .bib file
\usepackage{xcolor}
\usepackage{hyperref}
\hypersetup{
	colorlinks,
	citecolor={blue!70!black},
	linkcolor={blue!70!black},
	urlcolor={blue!70!black}}


% Set line spacing
\onehalfspacing % or use \doublespacing for double spacing

% Title information
\title{\textbf{\large{Working Title TBD\footnote{For useful comments we thank (Julien Senn), (Ernst Fehr), (Claus Kreiner), Jörg Oechssler, Dominika Langenmayr and participants at the 6th Workshop of the Swiss Network on Public Economics (Zurich), and the Workshop in Behavioral Public Economics and Political Economy (Vienna) }}}}
\author{\normalsize{Tim Kircali, Stefanie Stantcheva, Matthias Weber}}
\date{\normalsize{\today}}
\begin{document}
\maketitle
\pagenumbering{arabic}
%\begin{center}
%\textbf{Abstract} 
%\end{center}
%This study explores public preferences regarding tax enforcement and tax administration, focusing on attitudes towards the expansion of third-party reporting and the role of the tax authority in providing tax software and prepopulated tax returns. We conduct a survey experiment on a representative sample of the U.S.\ population. We vary the information that participants obtain to make causal inference on the role that information on tax evasion, inequality, and tax filing costs plays in views on comprehensive third-party reporting. Our results show that without additional information, about half of the respondents are in favor of comprehensive third-party reporting. This number increases up to almost three quarters for participants in our full information treatmet. Free tax software and the pre-population of tax returns by the tax authority are supported by a majority of respondents who have not received additional information and by a vast majority of respondents in our information treatments.


\section{Introduction}
Tax revenues are essential for governments to fulfill their responsibilities, such as providing security, infrastructure, education, health, and social protection. However, not all taxes that are due are actually paid---tax evasion poses a serious problem in many countries. Third-party reporting of tax-relevant information and international automatic exchange of information are very effective tools in reducing tax evasion \parencite{kleven2011unwilling, pomeranz2015no, slemrod2017does, boas2024taxing}. Moreover, comprehensive third-party reporting does not only help to enforce tax laws, but it can also improve the service that the tax authority can provide to citizens. In the U.S., taxpayers face high tax filing costs and are urged to purchase commercial tax software to file tax returns. On average, an individual taxpayer needs $13$ hours and about \$$250$ to file their taxes \parencite[][]{IRS2022}---when aggregated over all individual taxpayers in the U.S., this reflects a total cost of approximately \$$100$ billion per year. % Rough calculations: 160 million tax payers * (USD250 + 13 * USD25) = USD 92 billion OR 165 * (250 + 13 * 30) = 105600
In other countries, governments ease the burden of tax filing by offering free tax software and sometimes even pre-populated tax returns. Despite the benefits of third-party reporting and automatic exchange of information, the U.S.\ and some other countries are still reluctant to adopt these practices more widely.

We study preferences relating to these key issues in tax administration and tax enforcement on a sample that is representative of the U.S.\ population. We examine attitudes toward the reporting of tax-relevant data to the tax authority by third parties, toward international exchange of information, and toward services that the tax authority could provide, namely the free provision of tax software and the offering of pre-populated tax returns (the pre-population of tax returns requires sufficient data being reported to the tax authority by third parties). In simple words, we analyze people's preferences regarding questions such as the following. Should institutions and companies be obliged to report tax-relevant data to the tax authority? What type of data and which taxpayers should be included in such reporting? Should the tax authority reduce tax filing costs by providing free tax software? Should the tax authority offer pre-populated tax returns?

We leverage advances in survey methodology that allow us to uncover considerations and beliefs that drive the support for these practices. Therefore, we do not only ask about policy views but also ask questions relating to the factors and perceptions underlying their considerations and attitudes toward policy. In particular, we concentrate on respondents' considerations and beliefs about three main factors underlying policy support: (1) considerations of tax evasion with a focus on generated tax revenue, (2) considerations of tax evasion with a focus on fairness and inequality, and (3) considerations of tax filing costs (including the effort needed for tax filing). Additionally, we include questions to elicit their (policy-related) data privacy concerns, their perceived exposure to taxes, and questions to measure their prior knowledge about underlying factors such as the extent and distribution of tax evasion. 

%%Building on previous survey evidence we start our survey by asking questions about their demographic background, political leanings, government views; Factors that have proven to play a strong role in the formation of policy preferences \parencite{alesina2018intergenerational,stantcheva2021understanding}. 

To establish causal links between considerations and support for the policies, and to understand the role of misperceptions, we experimentally provide participants with information. We randomize participants to five groups, one control group receiving no information and four groups receiving information through informational video treatments. The videos provide information regarding the three analyzed main factors underlying policy support. (1) The evasion: revenue treatment provides information about the extent of tax evasion and the economic effects of information exchange on tax revenue.  (2) The evasion: inequality treatment discusses considerations related to fairness and inequality and provides information about the economic effects of information exchange on tax revenue. (3) The tax filing treatment discusses tax filing costs and provides information about tax software and prepopulated tax returns. And, (4) the full information treatment provides information on all three factors. The observed treatment differences in elicited considerations and policy views allow us (i) to uncover how people's considerations change when presented with specific information, (ii) to make causal inferences on which of the underlying considerations are driving policy views, and (iii) to explore how different considerations translate into policy views when weighed against one another.

%We try to understand the considerations driving the support for these practices and what role (missing) information plays in this support. In particular, we aim to understand the role of the following factors underlying policy support: (i) considerations of tax evasion with a focus on generated tax revenue, (ii) considerations of tax evasion with a focus on fairness and inequality, and (iii) considerations of tax filing costs (including the effort needed for tax filing). To establish causal links from these considerations and the related (missing) information to peoples support for the policies, we randomize participants to five treatments, one control treatment and four information treatments. In the four information treatments, participants see informational video treatments. In three treatments, the videos provide information regarding each of the three analyzed factors (evasion with a focus on revenue, evasion with a focus on fairness and inequality, and filing costs). The fourth treatment is a full information treatment providing information on all three factors. The observed treatment differences allow us to make causal inferences on which of the underlying considerations are primarily driving policy views. 

%In addition, previous research suggests that political leanings, government views, and the demographic background might play a role in the formation of the policy preferences. Therefore, we include analyses on how these shape the underlying considerations and the policy views. 


%\textbf{Preliminary Results.}
Our main findings can be summarized as follows. Considerations of tax evasion---both with a focus on generated tax revenue and on fairness and inequality---are strong predictors of support for comprehensive third-party reporting. These considerations often go hand in hand: people who expect third-party reporting to strongly reduce tax evasion also tend to believe that tax evasion is unevenly distributed, that comprehensive third-party reporting increases fairness, and that it is especially justified when limited to large businesses.

At the same time, people generally do not expect third-party reporting to have positive economic effects on themselves. Some subgroups express substantial data privacy concerns, and most respondents do not associate third-party reporting with potential improvements in filing procedures. Support for international automatic exchange of information is, on average, lower than for domestic third-party reporting. This difference is mainly driven by stronger data privacy concerns related to international data sharing.

People's support for tax software and pre-populated tax returns is generally very high. High-income individuals and those using tax preparers tend to be somewhat less supportive and prefer solutions provided by private companies. In contrast, low-income respondents and those filing taxes themselves show stronger support. Those in favor of pre-populated returns update their support for third-party reporting when informed that this is a key requirement for the pre-population of tax reports.

Turning to our causal analysis, we observe that the tax evasion treatment focusing on revenue increases support for third-party reporting and international automatic exchange of information to a similar extent. These shifts can be mainly explained by changes in the consideration of tax evasion, focusing on generated tax revenue. The inequality treatment, which focuses on fairness and inequality, shifts considerations and policy views in the same direction, but the treatment effect is stronger than in the revenue-focused treatment. A key finding is the significant increase in support for the international exchange of information. The main reasons for this are the significant changes in related fairness considerations and updated beliefs about the uneven distribution of offshore tax evasion. %This finding is consistent with previous findings about the strong role of fairness and inequality in people's preferences for economic policies. % If we wanted to keep the last sentence in, we would need to provide references.

The first two treatments have negligible effects on considerations about tax filing costs. These considerations change strongly in response to the tax filing treatment, which presents information about filing costs, tax software, and pre-populated returns. In addition to increasing support for tax software and pre-population, this treatment also has a positive effect on support for third-party reporting. First, this is a result of updated understanding about the key requirement of information exchange for the government to provide pre-populated tax returns. Second, it demonstrates how the anticipated personal benefits of tax software and pre-population motivate individuals to adjust their policy support for information exchange. 

%To preview our results, we find that without additional information (i.e., in our control treatment) about $54\%$ of respondents are in favor of comprehensive third-party reporting. When survey participants receive additional information on the consequences of such third-party reporting, the support increases to $62\%$ (tax evasion treatment), $64\%$ (inequality treatment), $69\%$ (tax filing treatment), or $72\%$ (in the full information treatment). For the international automatic exchange of information, the support is overall lower, but directions of treatments effect are as for third-party reporting: the support reaches from $39\%$ in the control treatment to $58\%$ in the full information treatment. Support for free tax software provided by the government is high in all treatments, with $78\%$ of respondents considering this valuable in the control treatment and up to $86\%$ in the full information treatment. Support for pre-populated tax returns is similarly high in all treatments, reaching from $76\%$ in the control treatment up to $88\%$ in the full information treatment.

%The role that fairness and inequality play in the support for third-party reporting can be seen by the fact that respondents are particularly in favor of third-party reporting if only data related to large businesses is included (in all treatments, but particularly so in the treatments with information on tax evasion). The considerations of tax filing costs play a very important role in people's decisions. The support for pre-populated tax returns is considerably higher than the support for third-party reporting in general, which might appear to be an inconsistency. However, when respondents are asked whether they would support comprehensive third-party reporting under the assumptions that the data will be used to provide pre-populated tax returns, the support for third-party reporting increases considerably over the general support in all treatments. The support for third-party reporting is also high if taxpayers could voluntarily choose to have their own data covered by third-party reporting (which would thus help to ease tax filing but not or hardly to reduce tax evasion).

%We observe some heterogeneity across the surveyed population. Politically left-leaning respondents are more supportive of all practices than right-leaning respondents. Younger respondents tend to be more supportive than older respondents and women more than men. We also find that income and level of education are positively associated with support of information sharing with the tax authority and of the tax authority's active provision of services to the taxpayers.

Our paper contributes to the literature in several ways. We provide the first measurement of support for third-party reporting and the international exchange of tax-related information, which are of crucial importance to reducing tax evasion. Furthermore, we provide novel evidence on the support for pre-populated tax returns and government-provided free tax software. We also explore the interplay of tax evasion considerations with compliance cost considerations (such as analyzing whether the support for third-party reporting increases when taxpayers are aware that it is a prerequisite for government-provided pre-populated tax returns, or if third-party reported data were only used to reduce tax compliance costs). By eliciting underlying beliefs and by experimentally providing policy-relevant information, we shed light on which considerations drive policy views and can draw causal inferences on which information influences policy views.



%\begin{center}
\textbf{Related Literature.}
%\end{center} 
Our paper relates to several strands of literature. First, this study contributes to the growing body of research that measures support for economic policies using survey experiment. A part of this literature analyzes preferences over taxation and taxation-related policies. \textcite{stantcheva2021understanding} explores beliefs, reasoning and views on income taxation and estate taxation. \textcite{cruces2013biased} and \textcite{kuziemko2015elastic} analyze inequality perceptions and support of redistributive policies, \textcite{alesina2018intergenerational} and \textcite{alesina2023immigration} analyze the effects of perceptions of intergenerational mobility and immigration on preferences for redistribution. Attitudes toward tax progressivity are examined in \textcite{tarroux2019value}. \textcite{liscow2022psychology} elicit preferences over the realization rule in capital gains taxation. \textcite{fisman2020americans} explore views toward a wealth tax in the U.S. Beyond taxation, survey experiments have, for instance, been used to analyze support for legalizing payments to kidney donors \parencite{elias2019paying} or views on climate policies \parencite{dechezlepretre2025fighting}. Much of the recent research in this field builds on advances in survey methodology, which facilitate uncovering perceptions, attitudes, and reasoning underlying policy views (see \textcite{haaland2023designing} for a discussion on designing information provision experiments and \textcite{stantcheva2023run} for a guide to run survey experiments). 

Second, this project is related to the literature on tax enforcement and tax evasion, especially to the literature on the relevance of information reporting. \textcite{slemrod2007cheating} documents that misreported income in the U.S.\ varies considerably by income type. While income from wages, salaries, interest and dividends (all of which are subject to third-party reporting) is rarely misreported, income from self-employment (typically self-reported) has estimated tax gaps above $50$\%. Similarly, \textcite{kleven2011unwilling} discover near-complete compliance and no deterrence effects for third-party reported income sources in Denmark. Relatedly, \textcite{pomeranz2015no} looks at VAT paper trails in Chile and finds that the effects of randomized deterrence messages are much lower in areas covered by VAT paper trails, suggesting that evasion is much lower in these areas. Other prior work investigates effects of information mechanisms such as consumer receipt collection \parencite{naritomi2019consumers}, the role of digital payment reporting forms like 1099-K \parencite{slemrod2017does}, and discusses the theoretical foundations for why third-party reporting improves compliance \parencite{kleven2016can}. Finally, \textcite{jensen2022employment} documents how modern tax systems arise over time and finds evidence that a rise in third-party reported income drives an expansion of the income tax base. Recent research also shows that the international automatic exchange of information (AEoI) reduces offshore tax evasion \parencite{boas2024taxing,baselgia2023compliance}.\footnote{A key finding of the analysis of offshore tax evasion is that it is much more prevalent among the highest income earners. While \textcite{johns2010distribution} reveal that the net misreporting percentage generally increases with income group but peaks at the $99$th percentile, \textcite{alstadsaeter2019tax} analyze leaked data and conclude that offshore tax evasion is highly concentrated among the super-rich. They estimate that the top $0.01$\% of households evade approximately $25$\% of their taxes. By contrast, tax evasion detected in stratified random tax audits is less than $5$ percent throughout the distribution. These results are consistent with those of \textcite{boas2024taxing}, who analyze offshore tax evasion in Denmark, and \textcite{guyton2021tax}, who study tax evasion at the top of the income distribution in the U.S. They show that previous estimates had failed to capture sophisticated evasion via offshore accounts and pass-through businesses, therefore underestimating top-$1$\% tax evasion.}

Third, this project also relates to the literature on tax compliance costs and the government's role in providing services to reduce them. In an early study on the tax compliance costs of individual taxpayers in the U.S., \textcite{blumenthal1992compliance} find that in 1989 taxpayers needed on average about $27$ hours and \$$66$ (1989 dollars) to file their taxes. In more recent work to understand how costly individuals perceive filing tax returns, \textcite{benzarti2020taxing} observes how taxpayers choose between itemizing deductions and claiming the standard deduction, finding that taxpayers forego large tax savings to avoid compliance costs. \textcite{hauck2024optional} shows that non-filing of tax declarations is particularly common at the lower end of the income distribution in Germany, which reduces the effectiveness of redistributive features of the tax system. \textcite{blesse2019people} present survey evidence eliciting the preferences of German taxpayers regarding the trade-off between simplifying the German tax code and reducing its progressivity. They discover that, while most taxpayers support simplifying the tax system, this support is partly driven by misconceptions about the trade-offs involved in tax complexity. \textcite{benzarti2024rising} ask U.S.\ respondents about the complexity of the tax code. They find that taxpayers perceive an increase in tax complexity and filing costs, that the vast majority of respondents consider tax filing to be tedious, and that almost all respondents would prefer a simpler tax system. Respondents would be willing to pay on average \$$130$ for a simpler tax system (despite the fact that the authors tell respondents that the simpler tax system might come with a higher tax liability due to fewer deduction possibilities). The majority of respondents is also supportive of pre-populated tax returns.



%\newpage
\printbibliography
\end{document}